\chapter{Oral Leukoplakia}

\section*{Definition}
Oral leukoplakia is a predominantly white, non-scrapable mucosal patch of the oral cavity that cannot be characterized clinically or histologically as any other definable disease, carrying a risk of malignant transformation estimated at 1--5\% over 5 years.

\section*{What Happens in Your Body}
Chronic mucosal irritation triggers hyperkeratosis (increased keratin production), acanthosis (epithelial thickening), and potentially dysplasia (disordered cellular maturation). Dysplastic changes progress from mild to severe dysplasia to carcinoma in situ and eventually invasive squamous cell carcinoma.

\section*{Causes}
\begin{itemize}
  \item \textbf{Tobacco use:} Most significant risk factor -- smoking, chewing tobacco, snuff, bidis
  \item \textbf{Alcohol consumption:} Synergistic effect with tobacco; risk ratio increases multiplicatively
  \item \textbf{Chronic mechanical irritation:} Ill-fitting dentures, sharp teeth, cheek biting
  \item \textbf{Human papillomavirus (HPV):} Subtypes 16 and 18 associated with some oral lesions
  \item \textbf{Actinic cheilitis:} Sun-induced leukoplakia of the lower lip
  \item \textbf{Immunosuppression:} HIV, post-transplant, or idiopathic CD4 lymphocytopenia
\end{itemize}

\section*{When to See a Doctor}
See your doctor if:
\begin{itemize}
  \item Any white oral mucosal patch that does not scrape off (distinguish from candidiasis)
  \item Lesion > 2 cm or located on the floor of mouth, ventral tongue, or soft palate (high-risk sites)
  \item Non-homogeneous (speckled, nodular, verrucous) or erythroleukoplakia (mixed red and white)
  \item Duration > 2 weeks without regression after removing suspected irritants
  \item Previous history of oral cancer or head and neck radiation
\end{itemize}

\section*{Related Symptoms}
\begin{itemize}
  \item White or gray-white patch on oral mucosa (buccal, gingival, palatal, lingual)
  \item Usually asymptomatic (painless)
  \item Rough, thickened, or fissured surface texture
  \item Burning sensation (atrophic or erosive subtypes)
  \item Erythroleukoplakia (mixed red and white) -- higher malignant risk
\end{itemize}
\textbf{Important:} Erythroleukoplakia (speckled leukoplakia) has a much higher malignant transformation risk (up to 25\%) than homogeneous leukoplakia and must be biopsied urgently.

\section*{Self-Care / Home Management}
\begin{itemize}
  \item Tobacco cessation in any form (counseling, NRT, varenicline, bupropion)
  \item Reduction or elimination of alcohol consumption
  \item Removal of mechanical irritants: denture adjustment, smoothing of sharp tooth edges
  \item Sun protection (zinc oxide lip balm) for actinic cheilitis
  \item No home treatment for the leukoplakia itself -- requires professional evaluation and biopsy
  \item Monthly self-examination with a mirror to monitor for changes
\end{itemize}

\section*{How Your Doctor Will Evaluate You}
\begin{itemize}
  \item Complete oral examination: note size, site, color, texture, and homogeneity of the lesion
  \item Toluidine blue staining (vital staining) to identify high-risk areas for biopsy
  \item Incisional biopsy of the most suspicious area (often the reddest, thickest, or most irregular portion)
  \item Excisional biopsy for small (< 1 cm) lesions
  \item Histopathologic grading: hyperkeratosis, dysplasia (mild/moderate/severe), carcinoma in situ
  \item HPV testing and p16 immunohistochemistry if indicated
  \item Annual surveillance for all people with a history of leukoplakia
\end{itemize}

\section*{Treatment Options}
\begin{itemize}
  \item \textbf{Mild dysplasia:} Observation with risk factor modification and close surveillance (every 6 months)
  \item \textbf{Moderate-to-severe dysplasia:} Excisional biopsy or surgical laser ablation (CO2 or Nd:YAG)
  \item \textbf{Medical therapy:} Topical retinoids (tretinoin gel 0.05\%) may cause regression but have high recurrence after discontinuation
  \item \textbf{Photodynamic therapy:} Alternative for diffuse or multifocal disease
  \item \textbf{No proven chemoprevention:} Vitamin A, beta-carotene, and retinoids show regression but no reduction in transformation risk
  \item \textbf{Actinic cheilitis:} Vermilionectomy or CO2 laser ablation of the lip
  \item \textbf{Long-term follow-up:} Every 6--12 months with photographic documentation for life
\end{itemize}

\clearpage
